\scanpage{323}\mspage{4\deleted{\unsure{6}}7}
% archive sheet number: 322

ce qui sera évident \deleted{\unsure{en}} grâce au résultat (donné plus \unsure{bas})
\[
 A=A(\unsure{\mathcal{F}|F}_{W})\Longleftrightarrow A(\mathrm{Fibf}_{W})
 \qquad\text{et de \={m} pour }B,C,
\]
qui signifie donc que \deleted{$A$}
\[
 \begin{aligned}
 \mathcal{F}_{W}\subset\mathrm{Fib}_{W}&\quad\text{implique}\quad\mathrm{Fibf}_{W}\subset\mathrm{Fib}_{W}\\
 &\hspace{4em}\text{(donc }\mathrm{Fibf}_{W}=\mathrm{Fib}_{W}),
 \end{aligned}
\]
or $\mathrm{Fib}_{W}$ étant stable par composition, il en \unsure{est} de même pour $\mathrm{Fibf}_{W}$, qed.

\uline{Corollaire} \unsure{4.4.6|4.46}. Supposons que $\mathcal{F}_{W_{\unsure{0|\bullet}}}$ soit stable par composition, ce qui est le cas notamment (\unsure{grâ à|grâce à} 4.4.5.2 appliqué à $W_0$ \unsure{ald|a.l.d.|au l de} $W$) si $W_0$ satisfait $\mathrm{Loc}(7)$. Alors les conditions $A(\mathcal{F}_{W_0}),B(\mathcal{F}_{W_0}),C(\mathcal{F}_{W_0})$ \deleted{sont équivalentes} (que nous notons aussi simplement par $A_{W_0},B_{W_0},C_{W_0}$) sont équivalentes. (\uline{NB} Si on ne fait pas l'hypothèse précédente sur $W_0$, on aura seulement $A_{W_0}\Longrightarrow B_{W_0}\Longrightarrow C_{W_0}$, en vertu de 4.4.2).

4.4.7. On a fait \unsure{le tour}, à présent de toutes les conditions \deleted{que} \inserted{\unsure{du|de} type ``fibration''} sur $W$, ou sur le couple $(W,W_0)$, qui m'ont paru utiles. Je vais les déployer ici dans un diagramme (plus ou moins) récapitulatif:
